\documentclass[10pt,a4paper]{amsart}
\usepackage[margin=19mm]{geometry}
\usepackage{amsmath,amssymb,mathtools}
\usepackage[scr=rsfs,cal=euler]{mathalfa}
\usepackage{xfrac}
\usepackage[unicode,hidelinks,bookmarks=false]{hyperref}
\newcommand{\ie}{i.e.\ }
\newcommand{\cf}{cf.\ }
\newcommand{\resp}{respectively\ }
\newcommand{\Subsection}{\subsection}
\providecommand{\nobreakdash}{\nobreak\hbox{-}}
\setlength{\emergencystretch}{2em}
\multlinegap0pt
\usepackage{amssymb}
\usepackage{enumerate}
\usepackage{bbm}
\usepackage[all]{xy}
\newtheorem{lem}{Lemma}
\title{Some Remarks on Super $M_{p}$-groups}
\author{Xiaoyou Chen, A. R. Moghaddamfar}
\date{Source: arXiv:2602.07172v1 (CC BY 4.0)}
\begin{document}
\maketitle

\noindent\emph{Source and licence.} Some Remarks on Super $M_{p}$-groups. Version arXiv:2602.07172v1 (CC BY 4.0), distributed by arXiv under the Creative Commons Attribution 4.0 International licence, http://creativecommons.org/licenses/by/4.0/, as recorded in the arXiv licence metadata for this exact version and shown on the abstract page https://arxiv.org/abs/2602.07172v1. Full licence notice: \texttt{licenses/arxiv-2602.07172-CC-BY-4.0.txt}.\par\smallskip

\noindent\textit{Editorial scope.} Counted unit: the article's lem lemma2.4 (source0.tex lines 266--271). The article's complete original proof of it is delivered with it (source0.tex lines 281--306, one \texttt{proof} environment of 26 lines, which the article places away from the statement). No mathematics is written, repaired or re-derived: the statement and the proof are the article's own lines, in the article's own notation, and no retained line is substituted or re-flowed.  Retained uncounted as the source-local context the proof needs: lines 251--254.  Nothing was omitted from any retained span, so the package carries no omission record.  No retained line contains a citation, so the wrapper carries no bibliography and no bibliography entry.  No retained line contains a figure environment, an image-inclusion command or drawing code, so no image is delivered and the package declares no asset dependency.  Editorial formatting only, in addition: an AMS \texttt{amsart} preamble that restates the article's own declarations and macros used by the retained lines, namely the environments \texttt{lem} on separate counters, together with the standard packages those lines need.  The retained environments print in this document as Lemma 1 in order of appearance, whereas the article prints the same items with the numbering of its own full document.  The complete native source of the article as distributed in the arXiv e-print for this exact version is bundled unchanged as \texttt{sources/194155/source0.tex}.
\begin{lem}\label{lemma2.3}
Let $\varphi$ be a monomial Brauer character of a group $G$, and let $K\leq G$ be a subgroup of $G$.
If $\varphi_{K}$ is irreducible, then $\varphi_{K}$ is also a monomial Brauer character.
\end{lem}

\begin{lem}\label{lemma2.4}
Let $\mathcal{G}$ be a class of groups that is closed under normal subgroups,
and let $\mathcal{M}$ be the class of $M_{p}$-groups in $\mathcal{G}$.
Suppose that every group in $\mathcal{M}$ is a super $M_{p}$-group.
Then, $\mathcal{M}$ is closed under normal subgroups.
\end{lem}

\begin{proof}[Proof of Lemma \ref{lemma2.4}]
Fix a group $G\in \mathcal{M}$, and let $N$ be a normal subgroup of $G$.
If $N=G$, the result is trivial. Therefore, we assume that $N<G$,
and there exists a maximal normal subgroup $M$ of $G$ such that $N\leq M<G$.
Since $G$ is an $M_{p}$-group it follows that $G$ is solvable.
Thus, $|G: M|=r$ for some prime $r$.
Notice that $M\in \mathcal{G}$ because $\mathcal{G}$ is closed under normal subgroups.


Let $\theta$ be an irreducible $p$-Brauer character of $M$. Define the inertia group of $\theta$ in $G$ as
$$T={\rm I}_{G}(\theta)=\{g\in G\mid \theta^{g}=\theta\}.$$ 
Since $|G: M|$ is a prime, we conclude that either $T=G$ or $T=M$.
By standard Clifford theory, it follows that either $\theta$ extends to $G$
or $\theta$ induces irreducibly to $G$.

If $\theta$ extends to $G$,
there exists a Brauer character $\varphi\in {\rm IBr}(G)$ such that $\varphi_{M}=\theta$.
Since $G$ is an $M_{p}$-group, we can apply Lemma \ref{lemma2.3} to conclude that $\theta$ is monomial.
If $\theta$ induces irreducibly to $G$, then $\theta^{G}\in {\rm IBr}(G)$.
Given that $G$ is a super $M_{p}$-group, we also deduce that $\theta$ is monomial.
Thus, we conclude that $M$ is an $M_{p}$-group, which implies that $M\in \mathcal{M}$.

Notice that $|M: N|<|G: N|$.
By induction on $|G: N|$ and the previous proof, we find that $N\in \mathcal{M}$.
Therefore, we can state that $\mathcal{M}$ is closed under normal subgroups.
\end{proof}

\end{document}
